% !TEX root = ../main.tex
\chapter{Introducció}\label{ch:introduccio}
\fancyhead[RE]{\slshape \textsf{Capítol \thechapter.\ Introducció}}
\fancyhead[LO]{\slshape \textsf{\nouppercase \rightmark}}

\section{Context i motivació}\label{sec:context}

Els instruments d'observació de la Terra produeixen imatges amb una resolució
espacial, espectral i temporal cada vegada més elevada. Aquest increment
millora la capacitat de caracteritzar vegetació, aigua, sòl, atmosfera o
infraestructures, però també augmenta el volum que cal emmagatzemar i
transmetre. En una missió espacial, la capacitat de l'enllaç de baixada no
creix necessàriament al mateix ritme que el sensor. Per tant, la compressió no
és únicament una etapa posterior de tractament de dades: condiciona quantes
imatges poden adquirir-se, conservar-se i enviar-se durant la vida de la
missió \cite{blanes2014tutorial,bioucas2013hyperspectral,denis2017disruptions}.

El problema s'accentua en nanosatèl·lits i altres plataformes amb recursos
limitats. La memòria, la potència elèctrica, la capacitat de càlcul i el temps
de contacte amb les estacions de terra són restriccions simultànies
\cite{selva2012cubesats,nasa2017cubesat}. Comprimir a terra no resol aquesta
situació, perquè les dades crues ja han hagut de travessar l'enllaç. La decisió
s'ha de prendre a bord, abans de transmetre.

Els estàndards CCSDS ofereixen solucions consolidades per a compressió sense
pèrdues, amb pèrdues i gairebé sense pèrdues
\cite{ccsds121,ccsds122,ccsds1221,ccsds123}. Paral·lelament, els còdecs basats
en xarxes neuronals han demostrat una gran capacitat per aprendre transformades
adaptades a les dades i obtenir bons compromisos entre la mida codificada i la
qualitat de reconstrucció
\cite{balle2017end,balle2018hyperprior,mijares2023reduced}. SORTENY parteix
d'aquesta segona línia: combina transformades espectrals i espacials apreses
per comprimir imatges multiespectrals \cite{mijares2025sorteny}.

\section{Antecedent propi i problema pendent}\label{sec:antecedent_intro}

Aquest treball continua un desenvolupament propi anterior que va adaptar el
compressor de referència de SORTENY, escrit en Python i TensorFlow, a una
implementació nativa en C \cite{anez2026treball}. Aquell desenvolupament va
exportar els pesos entrenats a fitxers consumibles des de C, va implementar els
operadors de la transformada d'anàlisi i va reduir l'ús de memòria mitjançant la
reutilització de buffers. També va introduir paral·lelisme i va executar el
compressor en una Raspberry Pi 3B+.

En una imatge de prova, la versió C va reduir la memòria màxima aproximadament
un factor $9,2\times$ i va accelerar l'execució $1,47\times$ respecte de la
referència Python. Aquests resultats demostraven que la inferència del
compressor podia traslladar-se a una plataforma amb menys d'un gigabyte de
memòria, però no constituïen encara un còdec complet.

Partint d'aquell resultat, quedaven cinc problemes oberts: completar la
reconstrucció en C; fer que el fitxer generat pel compressor inclogués tot el
que el descompressor necessita per reconstruir la imatge, sense dades auxiliars
externes; permetre qualitats diferents dins d'una mateixa imatge; decidir
automàticament quines regions cal protegir; i mesurar tant la mida real de les
dades codificades com el cost computacional del sistema complet. Aquests
problemes defineixen el punt de partida de la memòria actual.

\section{Problema i preguntes de recerca}\label{sec:problema}

Un còdec de qualitat uniforme assigna el mateix compromís de distorsió a tota
la imatge. Aquesta decisió és senzilla i reproduïble, però ignora que una
escena pot contenir zones amb valors operatius diferents. En una aplicació de
seguiment de vegetació pot interessar conservar millor les zones vegetades i
acceptar més distorsió a la resta. En una adquisició afectada per núvols pot
ser preferible dedicar més dades a les zones no cobertes per núvols, on la
superfície continua sent observable.

La selecció manual de regions és útil com a instrument experimental, però no
és una solució operativa general. Un satèl·lit no disposa d'un operador que
inspeccioni cada imatge abans de comprimir-la. Per tant, el sistema ha de
resoldre automàticament dues decisions diferents: identificar quines regions
són d'interès i determinar quina qualitat s'assigna a aquestes regions i a la
resta de la imatge. Separar les dues decisions permet canviar el criteri
d'interès sense redissenyar el còdec.

La recerca s'articula al voltant de quatre preguntes:

\begin{enumerate}[label=\textbf{P\arabic*.},leftmargin=*]
  \item Es pot completar i validar en C el procés de compressió i reconstrucció
  de SORTENY?
  \item Es poden identificar automàticament regions útils sense intervenció
  manual per a cada imatge?
  \item Es pot reduir la mida codificada redistribuint la qualitat sense
  perjudicar la regió que es vol protegir?
  \item És assumible el cost d'aquesta decisió i del còdec complet en una
  plataforma amb recursos limitats?
\end{enumerate}

Les hipòtesis mesurables associades a aquestes preguntes es formulen al
Capítol~\ref{ch:metodologia}, una vegada definits el sistema, les magnituds i
els criteris experimentals.

\section{Objectius}\label{sec:objectius}

L'objectiu general és completar i avaluar una extensió en C de SORTENY capaç
d'assignar automàticament qualitats diferents a regions d'una imatge
multiespectral i d'executar-se en maquinari amb recursos limitats.

Els objectius específics són:

\begin{enumerate}[label=\textbf{O\arabic*.},leftmargin=*]
  \item Completar i validar el compressor i el descompressor C respecte del
  comportament esperat del model de referència.
  \item Automatitzar la identificació de regions d'interès i l'assignació
  espacial de qualitat, sense requerir una selecció manual per imatge.
  \item Mesurar el compromís entre mida codificada i qualitat de reconstrucció
  respecte d'una compressió uniforme.
  \item Optimitzar les operacions C de més cost, conservar el comportament
  funcional i mesurar temps i memòria en una Raspberry Pi.
  \item Determinar en quines condicions l'estalvi de transmissió compensa el
  càlcul addicional i identificar les limitacions que encara impedeixen una
  adopció operativa completa.
\end{enumerate}

\section{Abast del treball}\label{sec:abast}

L'avaluació utilitza imatges multiespectrals de vuit bandes i una
implementació C executada tant en un entorn local com en una Raspberry Pi 3B+.
El treball no entrena un nou model neuronal: parteix dels pesos existents i es
centra a completar el còdec, controlar espacialment la qualitat i validar-ne
el comportament. La mida codificada final es mesura amb el codificador de la
implementació de referència; el flux C actual s'avalua separadament i encara no
incorpora aquesta etapa final de codificació.

La Raspberry Pi s'utilitza com una plataforma representativa de baix recurs,
útil per mesurar memòria, paral·lelisme i colls d'ampolla. L'existència de
plataformes a bord amb tecnologies COTS i prestacions comparables, com C3SatP a
la missió Menut, justifica l'interès de la comparació
\cite{schwaller2017a53,serra2024c3satp}. No obstant això, aquesta prova no
constitueix una qualificació de vol: no avalua radiació, consum energètic amb
instrumentació externa, fiabilitat de missió ni temps real per a un sensor
concret.

\section{Metodologia i organització}\label{sec:metodologia_organitzacio}

El desenvolupament s'ha realitzat de manera incremental. Abans de cada canvi
s'ha fixat una referència; després s'ha modificat un component aïllat i s'han
comparat els artefactes resultants. Una modificació només s'ha acceptat quan
les proves funcionals, les reconstruccions i les mesures associades han
confirmat el comportament esperat. Les proves de major cost s'han reservat per
als candidats que ja havien superat comprovacions més curtes.

El Capítol~\ref{ch:estat_art} presenta les fonts, els estàndards i les
alternatives considerades. El Capítol~\ref{ch:fonaments} introdueix els
fonaments de compressió i qualitat necessaris. Els Capítols~\ref{ch:pipeline}
i~\ref{ch:control} descriuen l'evolució del pipeline C i el control espacial.
El Capítol~\ref{ch:metodologia} formalitza les hipòtesis i el disseny
experimental. Els Capítols~\ref{ch:resultats} i~\ref{ch:bord} presenten els
resultats de qualitat, compressió i execució en hardware limitat. Finalment,
el Capítol~\ref{ch:conclusions} respon les preguntes de recerca, comprova els
objectius i delimita el treball futur.
